% from ega1-5-fr.tex
\oldpage[I]{131}
$H^1(X,\mathscr{J})=0$. Or,
$H^1(X,\mathscr{J})=H^1(X_0,\mathscr{J})$ et nous avons vu que
$\mathscr{J}$ est un $\mathscr{O}_{X_0}$-Module quasi-cohérent. Notre
assertion résultera donc du

\begin{lemma}[5.1.9.2]
\label{I.5.1.9.2-fr}
Si $Y$ est un schéma affine et $\mathscr{F}$ un $\mathscr{O}_Y$-Module
quasi-cohérent, on a $H^1(Y,\mathscr{F})=0$.
\end{lemma}

Ce lemme sera démontré au chap. III, § 1, comme conséquence du théorème plus
général que $H^i(Y,\mathscr{F})=0$ pour tout $i>0$. Pour en donner une
démonstration indépendante, observons que $H^1(Y,\mathscr{F})$ s'identifie au
module $\operatorname{Ext}^1_{\mathscr{O}_Y}
(Y;\mathscr{O}_Y,\mathscr{F})$ des classes d'extensions du
$\mathscr{O}_Y$-Module $\mathscr{O}_Y$ par le $\mathscr{O}_Y$-Module
$\mathscr{F}$ (T, 4.2.3)~; tout revient donc à prouver qu'une telle extension
$\mathscr{G}$ est triviale. Or, pour tout $y\in Y$, il y a un voisinage $V$
de $y$ dans $Y$ tel que $\mathscr{G}|V$ soit isomorphe à
$\mathscr{F}|Y\oplus\mathscr{O}_Y|V$
(\hyperref[0.5.4.9-fr]{0, 5.4.9})~; on en conclut que $\mathscr{G}$ est un
$\mathscr{O}_Y$-Module quasi-cohérent. Si $A$ est l'anneau de $Y$, on a donc
$\mathscr{F}=\widetilde{M}$, $\mathscr{G}=\widetilde{N}$, où $M$ et $N$ sont
des $A$-modules, et par hypothèse $N$ est une extension du $A$-module $A$ par
le $A$-module $M$ (\hyperref[I.1.3.11-fr]{1.3.11}). Comme cette extension est
nécessairement triviale, le lemme est démontré, et par suite aussi
(\hyperref[I.5.1.9-fr]{5.1.9}).

\begin{corollary}[5.1.10]
\label{I.5.1.10-fr}
Soit $X$ un préschéma tel que $\mathscr{N}_X$ soit nilpotent. Pour que $X$
soit un schéma affine, il faut et il suffit que $X_{\mathrm{red}}$ le soit.
\end{corollary}

\subsection{Existence d'un sous-préschéma d'espace sous-jacent donné.}
\label{subsection:I.5.2-fr}

\begin{proposition}[5.2.1]
\label{I.5.2.1-fr}
Pour tout sous-espace localement fermé $Y$ de l'espace sous-jacent à un
préschéma $X$, il existe un sous-préschéma réduit et un seul de $X$ ayant $Y$
pour espace sous-jacent.
\end{proposition}

L'unicité résultant de (\hyperref[I.5.1.2-fr]{5.1.2}), il reste à prouver
l'existence du sous-préschéma en question.

Si $X$ est affine d'anneau $A$, et $Y$ fermé dans $X$, la proposition est
immédiate~: $\mathfrak{j}(Y)$ est le plus grand idéal
$\mathfrak{a}\subset A$ tel que $V(\mathfrak{a})=Y$, et il est égal à sa
racine (\hyperref[I.1.1.4-fr]{1.1.4 (i)}), donc
$A/\mathfrak{j}(Y)$ est un anneau réduit.

Dans le cas général, pour tout ouvert affine $U\subset X$ tel que $U\cap Y$
soit fermé dans $U$, considérons le sous-préschéma fermé $Y_U$ de $U$ défini
par le faisceau d'idéaux associé à l'idéal $\mathfrak{j}(U\cap Y)$ de $A(U)$,
qui est réduit. Montrons que, si $V$ est un ouvert affine de $X$ contenu dans
$U$, $Y_V$ est \emph{induit} par $Y_U$ sur $V\cap Y$~; or, ce préschéma
induit est un sous-préschéma fermé de $V$ qui est réduit et a $V\cap Y$ comme
espace sous-jacent~; l'unicité de $Y_V$ entraîne donc notre assertion.

\begin{proposition}[5.2.2]
\label{I.5.2.2-fr}
Soient $X$ un préschéma réduit, $f:X\to Y$ un morphisme, $Z$ un
sous-préschéma fermé de $Y$ tel que $f(X)\subset Z$~; alors $f$ se factorise
en $X\xrightarrow{g}Z\xrightarrow{j}Y$, où $j$ est le morphisme d'injection.
\end{proposition}

Il résulte de l'hypothèse que le sous-préschéma fermé $f^{-1}(Z)$ de $X$ a
pour espace sous-jacent $X$ tout entier
(\hyperref[I.4.4.1-fr]{4.4.1})~; comme $X$ est réduit, ce sous-préschéma fermé
coïncide avec $X$ (\hyperref[I.5.1.2-fr]{5.1.2}), et la proposition résulte
donc de (\hyperref[I.4.4.1-fr]{4.4.1}).

\begin{corollary}[5.2.3]
\label{I.5.2.3-fr}
Soit $X$ un sous-préschéma réduit d'un préschéma $Y$~; si $Z$ est le
sous-préschéma fermé réduit de $Y$ ayant pour espace sous-jacent
$\overline{X}$, $X$ est un sous-préschéma induit sur un ouvert de $Z$.
\end{corollary}

